\section{\bt{Related Work}{相关工作}}
\label{sec:related}
\bi{Section~\ref{sec:alternatives} compares \system\ with agent memories, semantic layers, and caches (Table~\ref{tab:capabilities}). This section covers three further areas.}
{第~\ref{sec:alternatives} 节已将 \system\ 与智能体记忆、语义层和缓存进行比较（表~\ref{tab:capabilities}）。本节讨论另外三个方向。}
\paragraph{\bt{Feature engineering and temporal reuse}{特征工程与时间约束下的复用}}
\bi{Automated feature engineering and data preparation reduce the effort of building predictive inputs. CAAFE generates features from semantic information~\cite{hollmann2023caafe}; DS-Agent adapts successful data-science cases~\cite{guo2024dsagent}; FeatPilot augments a training table with features from related tables~\cite{liang2025featpilot}; DeepPrep builds preparation pipelines from execution feedback~\cite{fan2026deepprep}; KnowTrans transfers knowledge across data-preparation tasks~\cite{ge2025knowtrans}. RelBench evaluates relational prediction under temporal protocols~\cite{robinson2024relbench}, and feature stores such as Feast retrieve historical feature values with point-in-time joins~\cite{feast2026pointintime}. \system{} retains feature computations learned in agent sessions and reuses them across tasks and prediction cutoffs, compiling each definition with per-instance time bounds.}
{自动特征工程与数据准备降低了构造预测输入的工作量。CAAFE 利用语义信息生成特征~\cite{hollmann2023caafe}；DS-Agent 适配已有的成功数据科学案例~\cite{guo2024dsagent}；FeatPilot 从相关表中为训练表补充特征~\cite{liang2025featpilot}；DeepPrep 利用执行反馈构建数据准备流水线~\cite{fan2026deepprep}；KnowTrans 在数据准备任务之间迁移知识~\cite{ge2025knowtrans}。RelBench 以时间评估协议评测关系预测~\cite{robinson2024relbench}，Feast 等特征存储系统通过按时间点连接检索历史特征值~\cite{feast2026pointintime}。\system{} 保留智能体会话中学得的特征计算，并在后续任务和预测截点之间复用，以逐实例时间界限编译每个定义。}
\paragraph{\bt{Query reuse and result maintenance}{查询复用与结果维护}}
\bi{Multi-query optimization and shared execution reuse overlapping computation~\cite{sellis1988mqo,giannikis2012shareddb}, recycling retains intermediate results~\cite{ivanova2010recycling}, Ferdinand coordinates query-result caching and invalidation~\cite{garrod2008ferdinand}, and incremental view maintenance updates materialized results as base data change~\cite{blakeley1986efficiently,gupta1995maintenance}. These techniques keep results consistent with a stored query; \system{} checks whether the query still computes the intended metric, and the SQL it compiles still benefits from them.}
{多查询优化与共享执行复用重叠计算~\cite{sellis1988mqo,giannikis2012shareddb}，中间结果回收保留中间结果~\cite{ivanova2010recycling}，Ferdinand 协调查询结果缓存与失效处理~\cite{garrod2008ferdinand}，增量视图维护随基表变化更新物化结果~\cite{blakeley1986efficiently,gupta1995maintenance}。这些技术使结果与所存查询保持一致；\system{} 则检查该查询是否仍计算预期的指标，其编译出的 SQL 仍可受益于这些技术。}
\paragraph{\bt{Data conditions and validation}{数据条件与验证}}
\bi{Summarizability characterizes the conditions for correct aggregation~\cite{lenz1997summarizability}, and Huang et al.\ show how join fan-out double-counts aggregates in semantic layers~\cite{huang2023aggregation}. Constraint maintenance derives production rules that detect violations and restore consistency~\cite{ceri1990constraint}. Deequ supports declarative constraints, constraint suggestion, and incremental validation with shared scans~\cite{schelter2018deequ}; dbt data tests assert uniqueness and referential relationships over warehouse tables~\cite{dbt2026tests}. These tools check constraints declared or suggested for tables; \system{} derives conditions from each learned definition, records which definitions depend on each condition, and shares check results across definitions and agents.}
{汇总性刻画正确聚合所需的条件~\cite{lenz1997summarizability}，Huang 等人说明了语义层中连接扇出如何重复计算聚合值~\cite{huang2023aggregation}。约束维护派生产生式规则，用于检测违反并恢复一致性~\cite{ceri1990constraint}。Deequ 支持声明式约束、约束建议和基于共享扫描的增量验证~\cite{schelter2018deequ}；dbt data tests 在数据仓库表上断言唯一性与参照关系~\cite{dbt2026tests}。这些工具检查为表声明或建议的约束；\system{} 则从每个学得的定义中派生条件，记录每个条件支撑哪些定义，并在定义与智能体之间共享检查结果。}
